\chapter{Result-by-result derivation map}\label{coverage:chapter}

This map covers all 53 records in the existing result catalogue. Some describe the same mathematics at different locations, and some describe figures, estimation procedures or teaching resources. They are not 53 distinct theorems. Each link goes to the appropriate proof or calculation.

\section{Analytical results, examples and resources}

\small\begin{longtable}{@{}p{0.60\textwidth}p{0.32\textwidth}@{}}

\toprule Result or resource & Derivation or calculation\\\midrule\endhead

Main Figure 1: Construct conditional expected fitness before taking logs\newline{\footnotesize\itshape Worked or numerical calculation} & \hyperref[worked:main1]{p.~\pageref*{worked:main1}}; \hyperref[conditioning:environmental_average]{p.~\pageref*{conditioning:environmental_average}}\\[4pt]

Main Figure 2: Gaussian population projection identity\newline{\footnotesize\itshape Worked or numerical calculation} & \hyperref[gaussian:affine_projection]{p.~\pageref*{gaussian:affine_projection}}; \hyperref[worked:main2]{p.~\pageref*{worked:main2}}\\[4pt]

Main Figure 3: Proportional covariance contraction on a fixed surface\newline{\footnotesize\itshape Worked or numerical calculation} & \hyperref[geometry:scaling]{p.~\pageref*{geometry:scaling}}; \hyperref[worked:main3]{p.~\pageref*{worked:main3}}\\[4pt]

Main Figure 4: Scalar covariance trajectory with a separately changing gradient\newline{\footnotesize\itshape Worked or numerical calculation} & \hyperref[inheritance:scalar]{p.~\pageref*{inheritance:scalar}}; \hyperref[worked:main4]{p.~\pageref*{worked:main4}}\\[4pt]

Conditional mean, logarithm and reproductive weighting\newline{\footnotesize\itshape Analytical derivation} & \hyperref[conditioning:target]{p.~\pageref*{conditioning:target}}; \hyperref[conditioning:selection_weight]{p.~\pageref*{conditioning:selection_weight}}; \hyperref[conditioning:jensen]{p.~\pageref*{conditioning:jensen}}\\[4pt]

Gaussian environmental averaging changes breeding-value curvature\newline{\footnotesize\itshape Analytical derivation} & \hyperref[conditioning:environmental_average]{p.~\pageref*{conditioning:environmental_average}}; \hyperref[worked:main1]{p.~\pageref*{worked:main1}}\\[4pt]

Local quadratic approximation and the full nonquadratic curve\newline{\footnotesize\itshape Analytical derivation} & \hyperref[gaussian:taylor_derivation]{p.~\pageref*{gaussian:taylor_derivation}}; \hyperref[estimation:quartic_centring]{p.~\pageref*{estimation:quartic_centring}}\\[4pt]

Linear and centred-quadratic Gaussian moments\newline{\footnotesize\itshape Analytical derivation} & \hyperref[gaussian:moments]{p.~\pageref*{gaussian:moments}}; \hyperref[gaussian:variance_derivation]{p.~\pageref*{gaussian:variance_derivation}}\\[4pt]

Affine regression intercept, slope and R\textsuperscript{2}\newline{\footnotesize\itshape Analytical derivation} & \hyperref[gaussian:projection_section]{p.~\pageref*{gaussian:projection_section}}\\[4pt]

Noise reduces both R\textsuperscript{2} values but their ratio recovers the share\newline{\footnotesize\itshape Analytical derivation} & \hyperref[conditioning:log_gap_model]{p.~\pageref*{conditioning:log_gap_model}}; \hyperref[gaussian:noisy_response]{p.~\pageref*{gaussian:noisy_response}}\\[4pt]

Proportional covariance scaling and the omitted fraction\newline{\footnotesize\itshape Analytical derivation} & \hyperref[geometry:scaling]{p.~\pageref*{geometry:scaling}}; \hyperref[worked:main3]{p.~\pageref*{worked:main3}}\\[4pt]

A smaller covariance can reduce the linear share\newline{\footnotesize\itshape Explicit counterexample} & \hyperref[geometry:anisotropic_counterexample]{p.~\pageref*{geometry:anisotropic_counterexample}}\\[4pt]

Changing trait units leaves the component variances unchanged\newline{\footnotesize\itshape Analytical derivation} & \hyperref[geometry:coordinate_change]{p.~\pageref*{geometry:coordinate_change}}\\[4pt]

Whitening, curvature eigenvalues and the separate Rayleigh quotient\newline{\footnotesize\itshape Analytical derivation} & \hyperref[geometry:whitening]{p.~\pageref*{geometry:whitening}}; \hyperref[geometry:physical_orientation]{p.~\pageref*{geometry:physical_orientation}}; \hyperref[worked:rotation]{p.~\pageref*{worked:rotation}}\\[4pt]

Zero terms, undefined correlation and singular covariance support\newline{\footnotesize\itshape Analytical derivation} & \hyperref[gaussian:share]{p.~\pageref*{gaussian:share}}; \hyperref[gaussian:singular_support]{p.~\pageref*{gaussian:singular_support}}\\[4pt]

Asymmetry and non-Gaussian fourth moments can change the partition\newline{\footnotesize\itshape Analytical derivation} & \hyperref[gaussian:non_gaussian]{p.~\pageref*{gaussian:non_gaussian}}\\[4pt]

Exact selected-parent mean and covariance\newline{\footnotesize\itshape Analytical derivation} & \hyperref[selection:integrability]{p.~\pageref*{selection:integrability}}; \hyperref[selection:complete_square]{p.~\pageref*{selection:complete_square}}; \hyperref[selection:normaliser]{p.~\pageref*{selection:normaliser}}; \hyperref[selection:selected_moments]{p.~\pageref*{selection:selected_moments}}\\[4pt]

A large share does not ensure an accurate parental shift or perfect rankings\newline{\footnotesize\itshape Explicit counterexample} & \hyperref[selection:linear_comparison]{p.~\pageref*{selection:linear_comparison}}; \hyperref[selection:97_50_example]{p.~\pageref*{selection:97_50_example}}; \hyperref[worked:guide]{p.~\pageref*{worked:guide}}\\[4pt]

Absolute covariance-weighted curvature bounds relative displacement error\newline{\footnotesize\itshape Analytical derivation} & \hyperref[selection:error_bound_section]{p.~\pageref*{selection:error_bound_section}}; \hyperref[selection:eigenbasis_error]{p.~\pageref*{selection:eigenbasis_error}}\\[4pt]

Fixed allelic effects do not imply fixed genic covariance\newline{\footnotesize\itshape Analytical derivation} & \hyperref[inheritance:genic]{p.~\pageref*{inheritance:genic}}\\[4pt]

Fixed-baseline covariance interval and matrix fixed-point bounds\newline{\footnotesize\itshape Analytical derivation} & \hyperref[inheritance:matrix-order]{p.~\pageref*{inheritance:matrix-order}}; \hyperref[inheritance:matrix-fixed-point]{p.~\pageref*{inheritance:matrix-fixed-point}}\\[4pt]

Scalar fixed point, contraction and a transient increase\newline{\footnotesize\itshape Analytical derivation} & \hyperref[inheritance:scalar]{p.~\pageref*{inheritance:scalar}}; \hyperref[inheritance:scalar-boundaries]{p.~\pageref*{inheritance:scalar-boundaries}}\\[4pt]

No-restoration boundary and retained null-space variation\newline{\footnotesize\itshape Analytical derivation} & \hyperref[inheritance:rho-one]{p.~\pageref*{inheritance:rho-one}}\\[4pt]

Prediction shrinkage can mimic proportional biological contraction\newline{\footnotesize\itshape Analytical derivation} & \hyperref[estimation:conditioning]{p.~\pageref*{estimation:conditioning}}; \hyperref[estimation:shrinkage]{p.~\pageref*{estimation:shrinkage}}\\[4pt]

Convert compatible phenotype regression moments to breeding-value parameters\newline{\footnotesize\itshape Analytical derivation} & \hyperref[conversion:regression]{p.~\pageref*{conversion:regression}}; \hyperref[conversion:phenotype_moments]{p.~\pageref*{conversion:phenotype_moments}}; \hyperref[conversion:breeding_moments]{p.~\pageref*{conversion:breeding_moments}}; \hyperref[conversion:inversion]{p.~\pageref*{conversion:inversion}}\\[4pt]

Appendix S1: Gaussian moments and population projections\newline{\footnotesize\itshape Analytical derivation} & \hyperref[gaussian:chapter]{Chapter~\ref*{gaussian:chapter}}; \hyperref[geometry:chapter]{Chapter~\ref*{geometry:chapter}}\\[4pt]

Appendix S2: Gaussian selection, environmental averaging and displacement bound\newline{\footnotesize\itshape Analytical derivation} & \hyperref[conditioning:environmental_average]{p.~\pageref*{conditioning:environmental_average}}; \hyperref[selection:integrability]{p.~\pageref*{selection:integrability}}; \hyperref[selection:selected_moments]{p.~\pageref*{selection:selected_moments}}; \hyperref[selection:error_bound_section]{p.~\pageref*{selection:error_bound_section}}\\[4pt]

Appendix S3: Fixed-baseline inheritance and its limits\newline{\footnotesize\itshape Analytical derivation} & \hyperref[inheritance:chapter]{Chapter~\ref*{inheritance:chapter}}\\[4pt]

Appendix S4: Forward-time finite-population approximation\newline{\footnotesize\itshape Worked or numerical calculation} & \hyperref[worked:simulation]{p.~\pageref*{worked:simulation}}\\[4pt]

Appendix S5: Moment conversion from phenotype regressions\newline{\footnotesize\itshape Analytical derivation} & \hyperref[conversion:chapter]{Chapter~\ref*{conversion:chapter}}\\[4pt]

Appendix S6: Interpretation, assumptions and boundaries\newline{\footnotesize\itshape Interpretation} & \hyperref[coverage:questions]{p.~\pageref*{coverage:questions}}; \hyperref[concepts:chapter]{Chapter~\ref*{concepts:chapter}}\\[4pt]

Appendix S7: Visual one-trait and rotation examples\newline{\footnotesize\itshape Worked or numerical calculation} & \hyperref[worked:one_trait]{p.~\pageref*{worked:one_trait}}; \hyperref[worked:rotation]{p.~\pageref*{worked:rotation}}\\[4pt]

SLiM approximation check over 4,000 generation records\newline{\footnotesize\itshape Worked or numerical calculation} & \hyperref[worked:simulation]{p.~\pageref*{worked:simulation}}\\[4pt]

One fixed curve with broad, narrow and peak-centred populations\newline{\footnotesize\itshape Worked or numerical calculation} & \hyperref[worked:one_trait]{p.~\pageref*{worked:one_trait}}\\[4pt]

Rotate the same covariance eigenvalues on a fixed surface\newline{\footnotesize\itshape Worked or numerical calculation} & \hyperref[worked:rotation]{p.~\pageref*{worked:rotation}}\\[4pt]

Average phenotype-dependent expected output at each breeding value\newline{\footnotesize\itshape Worked or numerical calculation} & \hyperref[conditioning:environmental_average]{p.~\pageref*{conditioning:environmental_average}}; \hyperref[worked:guide]{p.~\pageref*{worked:guide}}\\[4pt]

Compare the linear and centred quadratic contributions\newline{\footnotesize\itshape Worked or numerical calculation} & \hyperref[gaussian:variance_derivation]{p.~\pageref*{gaussian:variance_derivation}}; \hyperref[worked:guide]{p.~\pageref*{worked:guide}}\\[4pt]

Reproductive weighting gives a selected-parent distribution\newline{\footnotesize\itshape Worked or numerical calculation} & \hyperref[selection:linear_comparison]{p.~\pageref*{selection:linear_comparison}}; \hyperref[selection:97_50_example]{p.~\pageref*{selection:97_50_example}}; \hyperref[worked:guide]{p.~\pageref*{worked:guide}}\\[4pt]

Iterate the scalar covariance closure\newline{\footnotesize\itshape Worked or numerical calculation} & \hyperref[inheritance:guide-fixed-point]{p.~\pageref*{inheritance:guide-fixed-point}}; \hyperref[worked:guide]{p.~\pageref*{worked:guide}}\\[4pt]

From phenotype moments to log conditional-fitness parameters\newline{\footnotesize\itshape Worked or numerical calculation} & \hyperref[conversion:scalar_example]{p.~\pageref*{conversion:scalar_example}}; \hyperref[worked:guide]{p.~\pageref*{worked:guide}}\\[4pt]

Every numerical checkpoint and the multivariate extension\newline{\footnotesize\itshape Worked or numerical calculation} & \hyperref[worked:guide]{p.~\pageref*{worked:guide}}; \hyperref[gaussian:chapter]{Chapter~\ref*{gaussian:chapter}}; \hyperref[selection:chapter]{Chapter~\ref*{selection:chapter}}; \hyperref[inheritance:chapter]{Chapter~\ref*{inheritance:chapter}}; \hyperref[conversion:chapter]{Chapter~\ref*{conversion:chapter}}\\[4pt]

Two-stage estimation from repeated measurements and offspring counts\newline{\footnotesize\itshape Statistical procedure} & \hyperref[estimation:calibration]{p.~\pageref*{estimation:calibration}}; \hyperref[estimation:likelihood]{p.~\pageref*{estimation:likelihood}}; \hyperref[estimation:components]{p.~\pageref*{estimation:components}}; \hyperref[worked:estimation]{p.~\pageref*{worked:estimation}}\\[4pt]

Joint refitting and model-conditional uncertainty\newline{\footnotesize\itshape Statistical procedure} & \hyperref[estimation:bootstrap]{p.~\pageref*{estimation:bootstrap}}; \hyperref[worked:estimation]{p.~\pageref*{worked:estimation}}\\[4pt]

Latent integration, recovery and explicit numerical failures\newline{\footnotesize\itshape Statistical procedure} & \hyperref[estimation:quadrature]{p.~\pageref*{estimation:quadrature}}; \hyperref[estimation:score]{p.~\pageref*{estimation:score}}; \hyperref[estimation:integration_check]{p.~\pageref*{estimation:integration_check}}\\[4pt]

Model shape, diagnostics and the local-versus-full distinction\newline{\footnotesize\itshape Analytical derivation} & \hyperref[estimation:aic]{p.~\pageref*{estimation:aic}}; \hyperref[estimation:quartic_centring]{p.~\pageref*{estimation:quartic_centring}}; \hyperref[estimation:quartic]{p.~\pageref*{estimation:quartic}}; \hyperref[worked:estimation]{p.~\pageref*{worked:estimation}}\\[4pt]

Observed counts and latent conditional relationship\newline{\footnotesize\itshape Statistical procedure} & \hyperref[estimation:components]{p.~\pageref*{estimation:components}}; \hyperref[estimation:predictions]{p.~\pageref*{estimation:predictions}}; \hyperref[worked:estimation]{p.~\pageref*{worked:estimation}}\\[4pt]

Component estimates, naive comparison and conditional interval\newline{\footnotesize\itshape Statistical procedure} & \hyperref[estimation:bootstrap]{p.~\pageref*{estimation:bootstrap}}; \hyperref[worked:estimation]{p.~\pageref*{worked:estimation}}\\[4pt]

Observable-bin predictions and zero-count checks\newline{\footnotesize\itshape Statistical procedure} & \hyperref[estimation:predictions]{p.~\pageref*{estimation:predictions}}; \hyperref[estimation:bin_variance]{p.~\pageref*{estimation:bin_variance}}; \hyperref[estimation:quartic]{p.~\pageref*{estimation:quartic}}\\[4pt]

Runnable scientific source, reference data and result checks\newline{\footnotesize\itshape Reproduction resource} & \hyperref[worked:commands]{p.~\pageref*{worked:commands}}; \hyperref[coverage:commands]{p.~\pageref*{coverage:commands}}\\[4pt]

Observation means, population-fitted lines and parental shift\newline{\footnotesize\itshape Worked or numerical calculation} & \hyperref[conditioning:chapter]{Chapter~\ref*{conditioning:chapter}}; \hyperref[worked:one_trait]{p.~\pageref*{worked:one_trait}}; \hyperref[worked:guide]{p.~\pageref*{worked:guide}}\\[4pt]

Rotate, narrow or move a Gaussian population on one fixed surface\newline{\footnotesize\itshape Worked or numerical calculation} & \hyperref[geometry:chapter]{Chapter~\ref*{geometry:chapter}}; \hyperref[worked:rotation]{p.~\pageref*{worked:rotation}}; \hyperref[worked:comparison]{p.~\pageref*{worked:comparison}}\\[4pt]

Frequency, expected-output and recorded-count weights\newline{\footnotesize\itshape Worked or numerical calculation} & \hyperref[conditioning:selection_weight]{p.~\pageref*{conditioning:selection_weight}}; \hyperref[worked:parental]{p.~\pageref*{worked:parental}}\\[4pt]

Change measurement uncertainty, sample size and model shape\newline{\footnotesize\itshape Reproduction resource} & \hyperref[estimation:likelihood]{p.~\pageref*{estimation:likelihood}}; \hyperref[estimation:predictions]{p.~\pageref*{estimation:predictions}}; \hyperref[estimation:bootstrap]{p.~\pageref*{estimation:bootstrap}}; \hyperref[worked:commands]{p.~\pageref*{worked:commands}}\\[4pt]

\bottomrule\end{longtable}\normalsize

\section{Equation-by-equation map}\label{coverage:equations}

The following 106 source-qualified equation identifiers connect the main account and supplementary derivations to this companion. These are stable source labels, not printed equation numbers. The same label can occur in different sources. Definitions and model assumptions are identified separately from proofs and worked calculations. The source versions inspected on 2026-11-13 are identified by hashes in \repofile{tutorials/derivations/equation_map.json}{the equation map}. The map contains identifiers and concepts; the source documents are not needed to build this companion.

\Needspace{8\baselineskip}

\subsection{Main text and mathematical appendices}

\small\begin{longtable}{@{}>{\raggedright\arraybackslash}p{0.28\textwidth}>{\raggedright\arraybackslash}p{0.47\textwidth}>{\raggedright\arraybackslash}p{0.16\textwidth}@{}}

\toprule Source label & Concept and treatment & Derivation\\\midrule\endhead

\codepath{eq:target} & Conditional expected reproduction and its logarithm\newline{\footnotesize\itshape definition} & \hyperref[conditioning:target]{p.~\pageref*{conditioning:target}}; \hyperref[conditioning:log_target]{p.~\pageref*{conditioning:log_target}}\\[4pt]

\codepath{eq:selection_weight} & Normalised reproductive weighting of parental frequencies\newline{\footnotesize\itshape proof + definition} & \hyperref[conditioning:selection_weight]{p.~\pageref*{conditioning:selection_weight}}; \hyperref[conditioning:selected_density]{p.~\pageref*{conditioning:selected_density}}\\[4pt]

\codepath{eq:log_ratio} & Log differences as proportional changes in conditional mean output\newline{\footnotesize\itshape proof} & \hyperref[conditioning:log_ratio]{p.~\pageref*{conditioning:log_ratio}}\\[4pt]

\codepath{eq:jensen} & Noncommutation of averaging and taking logarithms\newline{\footnotesize\itshape proof} & \hyperref[conditioning:jensen]{p.~\pageref*{conditioning:jensen}}; \hyperref[conditioning:jensen_inequality]{p.~\pageref*{conditioning:jensen_inequality}}\\[4pt]

\codepath{eq:taylor} & Local quadratic expansion of log conditional expected output\newline{\footnotesize\itshape definition + proof} & \hyperref[gaussian:quadratic_definition]{p.~\pageref*{gaussian:quadratic_definition}}; \hyperref[gaussian:taylor_derivation]{p.~\pageref*{gaussian:taylor_derivation}}\\[4pt]

\codepath{eq:centered_terms} & Mean-centred quadratic component\newline{\footnotesize\itshape definition + proof} & \hyperref[gaussian:terms]{p.~\pageref*{gaussian:terms}}; \hyperref[gaussian:quadratic_mean]{p.~\pageref*{gaussian:quadratic_mean}}\\[4pt]

\codepath{eq:vlin} & Variance of the linear component\newline{\footnotesize\itshape proof} & \hyperref[gaussian:linear_variance]{p.~\pageref*{gaussian:linear_variance}}\\[4pt]

\codepath{eq:vquad} & Gaussian variance of the quadratic component\newline{\footnotesize\itshape proof} & \hyperref[gaussian:four_index_variance]{p.~\pageref*{gaussian:four_index_variance}}; \hyperref[gaussian:quadratic_variance]{p.~\pageref*{gaussian:quadratic_variance}}\\[4pt]

\codepath{eq:Ag} & First-order share of quadratic log-target variance\newline{\footnotesize\itshape definition} & \hyperref[gaussian:share]{p.~\pageref*{gaussian:share}}; \hyperref[gaussian:variance_partition]{p.~\pageref*{gaussian:variance_partition}}\\[4pt]

\codepath{eq:affine_projection} & Population affine projection and curvature-adjusted intercept\newline{\footnotesize\itshape proof} & \hyperref[gaussian:affine_projection]{p.~\pageref*{gaussian:affine_projection}}; \hyperref[gaussian:fitted_affine]{p.~\pageref*{gaussian:fitted_affine}}\\[4pt]

\codepath{eq:r2_identity} & Regression and squared-correlation interpretations of the share\newline{\footnotesize\itshape proof} & \hyperref[gaussian:regression_share]{p.~\pageref*{gaussian:regression_share}}; \hyperref[gaussian:correlation_identity]{p.~\pageref*{gaussian:correlation_identity}}\\[4pt]

\codepath{eq:nested_r2} & Ratio of nested fit statistics and the complementary fraction\newline{\footnotesize\itshape proof} & \hyperref[gaussian:nested_r2]{p.~\pageref*{gaussian:nested_r2}}; \hyperref[gaussian:nested_ratio]{p.~\pageref*{gaussian:nested_ratio}}\\[4pt]

\codepath{eq:scaling} & Linear and quadratic components under proportional covariance scaling\newline{\footnotesize\itshape proof} & \hyperref[geometry:scaling_proposition]{p.~\pageref*{geometry:scaling_proposition}}; \hyperref[geometry:component_scaling]{p.~\pageref*{geometry:component_scaling}}\\[4pt]

\codepath{eq:whitening} & Gradient and Hessian in population-standardised coordinates\newline{\footnotesize\itshape definition + proof} & \hyperref[geometry:whitened_coefficients]{p.~\pageref*{geometry:whitened_coefficients}}; \hyperref[geometry:spectral_components]{p.~\pageref*{geometry:spectral_components}}\\[4pt]

\codepath{eq:log_quadratic_weight} & Exact quadratic reproductive-weight model\newline{\footnotesize\itshape definition} & \hyperref[selection:model]{p.~\pageref*{selection:model}}; \hyperref[selection:log_quadratic_weight]{p.~\pageref*{selection:log_quadratic_weight}}\\[4pt]

\codepath{eq:selected_mean} & Exact selected-parent mean displacement\newline{\footnotesize\itshape proof} & \hyperref[selection:selected_moments]{p.~\pageref*{selection:selected_moments}}; \hyperref[selection:moment_formulas]{p.~\pageref*{selection:moment_formulas}}\\[4pt]

\codepath{eq:selected_covariance} & Exact covariance under parental reproductive weights\newline{\footnotesize\itshape proof} & \hyperref[selection:selected_moments]{p.~\pageref*{selection:selected_moments}}; \hyperref[selection:moment_formulas]{p.~\pageref*{selection:moment_formulas}}\\[4pt]

\codepath{eq:counterexample} & Large variance share with a one-half displacement overestimate\newline{\footnotesize\itshape calculation} & \hyperref[selection:97_50_example]{p.~\pageref*{selection:97_50_example}}; \hyperref[worked:guide]{p.~\pageref*{worked:guide}}\\[4pt]

\codepath{eq:T} & Population-scaled stabilising curvature\newline{\footnotesize\itshape definition} & \hyperref[selection:T]{p.~\pageref*{selection:T}}; \hyperref[selection:precision_factorisation]{p.~\pageref*{selection:precision_factorisation}}\\[4pt]

\codepath{eq:response_bound} & Covariance-scaled relative displacement-error bound\newline{\footnotesize\itshape proof} & \hyperref[selection:relative_bound]{p.~\pageref*{selection:relative_bound}}; \hyperref[selection:full_bound]{p.~\pageref*{selection:full_bound}}\\[4pt]

\codepath{eq:fixedB_recursion} & Covariance update under fixed genic baseline and association retention\newline{\footnotesize\itshape definition + proof} & \hyperref[inheritance:retention-assumption]{p.~\pageref*{inheritance:retention-assumption}}; \hyperref[inheritance:recursion]{p.~\pageref*{inheritance:recursion}}\\[4pt]

\codepath{eq:genic_baseline} & Genic covariance from allele frequencies and fixed allelic effects\newline{\footnotesize\itshape proof} & \hyperref[inheritance:genic-formula]{p.~\pageref*{inheritance:genic-formula}}; \hyperref[inheritance:association-decomposition]{p.~\pageref*{inheritance:association-decomposition}}\\[4pt]

\codepath{eq:loewner_interval} & Invariant order interval for the covariance map\newline{\footnotesize\itshape definition + proof} & \hyperref[inheritance:interval]{p.~\pageref*{inheritance:interval}}; \hyperref[inheritance:invariance]{p.~\pageref*{inheritance:invariance}}\\[4pt]

\codepath{eq:fixedpoint_bounds} & Bounds on every positive covariance fixed point\newline{\footnotesize\itshape proof} & \hyperref[inheritance:existence]{p.~\pageref*{inheritance:existence}}; \hyperref[inheritance:fixed-bounds]{p.~\pageref*{inheritance:fixed-bounds}}\\[4pt]

\codepath{eq:scalar_recursion} & One-trait or isotropic covariance update\newline{\footnotesize\itshape definition + proof} & \hyperref[inheritance:scalar-map]{p.~\pageref*{inheritance:scalar-map}}; \hyperref[inheritance:scalar]{p.~\pageref*{inheritance:scalar}}\\[4pt]

\codepath{eq:scalar_root} & Unique positive scalar covariance fixed point\newline{\footnotesize\itshape proof} & \hyperref[inheritance:scalar-quadratic]{p.~\pageref*{inheritance:scalar-quadratic}}; \hyperref[inheritance:positive-root]{p.~\pageref*{inheritance:positive-root}}\\[4pt]

\codepath{eq:rho_one} & Precision accumulation with complete association retention\newline{\footnotesize\itshape proof} & \hyperref[inheritance:precision-iteration]{p.~\pageref*{inheritance:precision-iteration}}; \hyperref[inheritance:rho-one]{p.~\pageref*{inheritance:rho-one}}\\[4pt]

\codepath{eq:log_error_model} & Observation model with a constant conditional log gap\newline{\footnotesize\itshape definition + proof} & \hyperref[conditioning:constant_gap]{p.~\pageref*{conditioning:constant_gap}}; \hyperref[gaussian:noise_model]{p.~\pageref*{gaussian:noise_model}}\\[4pt]

\codepath{eq:ebv_inflation} & Share after substituting shrunken prediction covariance\newline{\footnotesize\itshape proof} & \hyperref[estimation:shrinkage_formula]{p.~\pageref*{estimation:shrinkage_formula}}; \hyperref[estimation:shrinkage_share]{p.~\pageref*{estimation:shrinkage_share}}\\[4pt]

\codepath{eq:LA_regression} & Raw relative-fitness phenotype regression and quadratic convention\newline{\footnotesize\itshape definition} & \hyperref[conversion:regression_model]{p.~\pageref*{conversion:regression_model}}; \hyperref[conversion:coefficient_convention]{p.~\pageref*{conversion:coefficient_convention}}\\[4pt]

\codepath{eq:LA_conversion} & Conversion of selected moments into log-quadratic parameters\newline{\footnotesize\itshape proof} & \hyperref[conversion:selected_formula]{p.~\pageref*{conversion:selected_formula}}; \hyperref[conversion:conversion_formula]{p.~\pageref*{conversion:conversion_formula}}\\[4pt]

\codepath{eq:fitness_model} & Phenotype-fitness model with independent Gaussian environment\newline{\footnotesize\itshape definition} & \hyperref[conditioning:phenotype_fitness]{p.~\pageref*{conditioning:phenotype_fitness}}; \hyperref[conditioning:environmental_average]{p.~\pageref*{conditioning:environmental_average}}\\[4pt]

\codepath{eq:environmental_average} & Gaussian environmental convolution of the fitness curve\newline{\footnotesize\itshape proof} & \hyperref[conditioning:environment_formula]{p.~\pageref*{conditioning:environment_formula}}; \hyperref[conditioning:environmental_average]{p.~\pageref*{conditioning:environmental_average}}\\[4pt]

\codepath{eq:precision_order} & Inverse-order covariance contraction under stabilising curvature\newline{\footnotesize\itshape proof} & \hyperref[inheritance:inverse-order]{p.~\pageref*{inheritance:inverse-order}}; \hyperref[inheritance:selection-order]{p.~\pageref*{inheritance:selection-order}}\\[4pt]

\codepath{eq:LA_moments} & Phenotype-regression normalisation and selected moments\newline{\footnotesize\itshape proof} & \hyperref[conversion:regression_moments]{p.~\pageref*{conversion:regression_moments}}; \hyperref[conversion:phenotype_moment_formula]{p.~\pageref*{conversion:phenotype_moment_formula}}\\[4pt]

\codepath{eq:LA_selected_mean} & Selected breeding-value mean from phenotype regression\newline{\footnotesize\itshape proof} & \hyperref[conversion:selected_breeding_moments]{p.~\pageref*{conversion:selected_breeding_moments}}; \hyperref[conversion:selected_formula]{p.~\pageref*{conversion:selected_formula}}\\[4pt]

\codepath{eq:LA_selected_cov} & Selected breeding-value covariance from phenotype regression\newline{\footnotesize\itshape proof} & \hyperref[conversion:selected_breeding_moments]{p.~\pageref*{conversion:selected_breeding_moments}}; \hyperref[conversion:selected_formula]{p.~\pageref*{conversion:selected_formula}}\\[4pt]

\bottomrule\end{longtable}\normalsize

\Needspace{8\baselineskip}

\subsection{Appendix S7: visual walkthrough}

\small\begin{longtable}{@{}>{\raggedright\arraybackslash}p{0.28\textwidth}>{\raggedright\arraybackslash}p{0.47\textwidth}>{\raggedright\arraybackslash}p{0.16\textwidth}@{}}

\toprule Source label & Concept and treatment & Derivation\\\midrule\endhead

\codepath{eq:visual_model} & Fixed one-trait curve and Gaussian populations\newline{\footnotesize\itshape definition} & \hyperref[worked:one_trait]{p.~\pageref*{worked:one_trait}}\\[4pt]

\codepath{eq:visual_variances} & Variance components for the three one-trait populations\newline{\footnotesize\itshape calculation + proof} & \hyperref[worked:one_trait]{p.~\pageref*{worked:one_trait}}; \hyperref[gaussian:variance_formulas]{p.~\pageref*{gaussian:variance_formulas}}\\[4pt]

\codepath{eq:visual_line} & Population-fitted line for the one-trait walkthrough\newline{\footnotesize\itshape calculation + proof} & \hyperref[worked:one_trait]{p.~\pageref*{worked:one_trait}}; \hyperref[gaussian:fitted_affine]{p.~\pageref*{gaussian:fitted_affine}}\\[4pt]

\codepath{eq:visual_rotation_surface} & Two-trait surface expressed along and across the gradient\newline{\footnotesize\itshape definition + calculation} & \hyperref[worked:rotation]{p.~\pageref*{worked:rotation}}; \hyperref[geometry:rotation_formula]{p.~\pageref*{geometry:rotation_formula}}\\[4pt]

\codepath{eq:visual_rotation_variances} & Full-population components under physical covariance rotation\newline{\footnotesize\itshape calculation + proof} & \hyperref[worked:rotation]{p.~\pageref*{worked:rotation}}; \hyperref[geometry:rotation_formula]{p.~\pageref*{geometry:rotation_formula}}\\[4pt]

\bottomrule\end{longtable}\normalsize

\Needspace{8\baselineskip}

\subsection{Supplementary File 1: mathematics guide}

\small\begin{longtable}{@{}>{\raggedright\arraybackslash}p{0.28\textwidth}>{\raggedright\arraybackslash}p{0.47\textwidth}>{\raggedright\arraybackslash}p{0.16\textwidth}@{}}

\toprule Source label & Concept and treatment & Derivation\\\midrule\endhead

\codepath{eq:guide-target} & Conditional expected reproductive output\newline{\footnotesize\itshape definition} & \hyperref[conditioning:target]{p.~\pageref*{conditioning:target}}\\[4pt]

\codepath{eq:guide-log-target} & Logarithm of conditional expected output\newline{\footnotesize\itshape definition} & \hyperref[conditioning:log_target]{p.~\pageref*{conditioning:log_target}}\\[4pt]

\codepath{eq:guide-jensen} & Averaging versus logging individual outputs\newline{\footnotesize\itshape proof} & \hyperref[conditioning:jensen]{p.~\pageref*{conditioning:jensen}}; \hyperref[conditioning:jensen_inequality]{p.~\pageref*{conditioning:jensen_inequality}}\\[4pt]

\codepath{eq:guide-selection-weight} & Normalised reproductive weighting of breeding values\newline{\footnotesize\itshape proof + definition} & \hyperref[conditioning:selection_weight]{p.~\pageref*{conditioning:selection_weight}}; \hyperref[conditioning:selected_density]{p.~\pageref*{conditioning:selected_density}}\\[4pt]

\codepath{eq:guide-phenotype-fitness} & The guide's Gaussian phenotype-fitness example\newline{\footnotesize\itshape definition + calculation} & \hyperref[conditioning:phenotype_fitness]{p.~\pageref*{conditioning:phenotype_fitness}}; \hyperref[worked:guide]{p.~\pageref*{worked:guide}}\\[4pt]

\codepath{eq:guide-environmental-average} & Conditional mean for the guide's environmental parameters\newline{\footnotesize\itshape proof + calculation} & \hyperref[conditioning:environment_formula]{p.~\pageref*{conditioning:environment_formula}}; \hyperref[worked:guide]{p.~\pageref*{worked:guide}}\\[4pt]

\codepath{eq:guide-quadratic-example} & Expansion of the guide's exact quadratic log mean\newline{\footnotesize\itshape calculation} & \hyperref[worked:guide]{p.~\pageref*{worked:guide}}; \hyperref[conditioning:environment_derivatives]{p.~\pageref*{conditioning:environment_derivatives}}\\[4pt]

\codepath{eq:guide-finite-slope} & Symmetric finite-difference estimate of slope\newline{\footnotesize\itshape calculation + proof} & \hyperref[worked:guide]{p.~\pageref*{worked:guide}}; \hyperref[gaussian:taylor_derivation]{p.~\pageref*{gaussian:taylor_derivation}}\\[4pt]

\codepath{eq:guide-finite-curvature} & Symmetric second difference for curvature\newline{\footnotesize\itshape calculation + proof} & \hyperref[worked:guide]{p.~\pageref*{worked:guide}}; \hyperref[gaussian:taylor_derivation]{p.~\pageref*{gaussian:taylor_derivation}}\\[4pt]

\codepath{eq:guide-local-description} & One-trait quadratic approximation\newline{\footnotesize\itshape definition + proof} & \hyperref[gaussian:quadratic_definition]{p.~\pageref*{gaussian:quadratic_definition}}; \hyperref[gaussian:taylor_derivation]{p.~\pageref*{gaussian:taylor_derivation}}\\[4pt]

\codepath{eq:guide-environmental-integral} & Integral over environmental deviations in the guide example\newline{\footnotesize\itshape definition + proof} & \hyperref[conditioning:environment_integral]{p.~\pageref*{conditioning:environment_integral}}; \hyperref[conditioning:environmental_average]{p.~\pageref*{conditioning:environmental_average}}\\[4pt]

\codepath{eq:guide-general-completing-square} & Completing the environmental square\newline{\footnotesize\itshape proof} & \hyperref[conditioning:environmental_average]{p.~\pageref*{conditioning:environmental_average}}\\[4pt]

\codepath{eq:guide-environmental-average-general} & General Gaussian environmental convolution\newline{\footnotesize\itshape proof} & \hyperref[conditioning:environment_formula]{p.~\pageref*{conditioning:environment_formula}}; \hyperref[conditioning:environmental_average]{p.~\pageref*{conditioning:environmental_average}}\\[4pt]

\codepath{eq:guide-normal-moments} & Centred Gaussian moments through order four\newline{\footnotesize\itshape proof} & \hyperref[gaussian:mgf_derivatives]{p.~\pageref*{gaussian:mgf_derivatives}}; \hyperref[gaussian:fourth_moment]{p.~\pageref*{gaussian:fourth_moment}}\\[4pt]

\codepath{eq:guide-centered-Q} & Centred scalar quadratic component\newline{\footnotesize\itshape definition + proof} & \hyperref[gaussian:terms]{p.~\pageref*{gaussian:terms}}; \hyperref[gaussian:quadratic_mean]{p.~\pageref*{gaussian:quadratic_mean}}\\[4pt]

\codepath{eq:guide-vlin-one} & Scalar linear variance\newline{\footnotesize\itshape proof} & \hyperref[gaussian:linear_variance]{p.~\pageref*{gaussian:linear_variance}}\\[4pt]

\codepath{eq:guide-vquad-one} & Scalar quadratic variance\newline{\footnotesize\itshape proof} & \hyperref[gaussian:quadratic_variance]{p.~\pageref*{gaussian:quadratic_variance}}; \hyperref[gaussian:fourth_moment]{p.~\pageref*{gaussian:fourth_moment}}\\[4pt]

\codepath{eq:guide-correlation} & Squared correlation between the linear term and full quadratic target\newline{\footnotesize\itshape proof} & \hyperref[gaussian:correlation_identity]{p.~\pageref*{gaussian:correlation_identity}}\\[4pt]

\codepath{eq:Ag} & First-order variance share\newline{\footnotesize\itshape definition} & \hyperref[gaussian:share]{p.~\pageref*{gaussian:share}}; \hyperref[gaussian:singular_support]{p.~\pageref*{gaussian:singular_support}}\\[4pt]

\codepath{eq:guide-projection} & Population affine-regression R-squared equals the share\newline{\footnotesize\itshape proof} & \hyperref[gaussian:regression_share]{p.~\pageref*{gaussian:regression_share}}; \hyperref[gaussian:affine_projection]{p.~\pageref*{gaussian:affine_projection}}\\[4pt]

\codepath{eq:guide-one-trait-slope} & Scalar least-squares slope equals the local slope under Gaussian quadratic conditions\newline{\footnotesize\itshape proof} & \hyperref[gaussian:normal_equations]{p.~\pageref*{gaussian:normal_equations}}; \hyperref[gaussian:affine_projection]{p.~\pageref*{gaussian:affine_projection}}\\[4pt]

\codepath{eq:guide-nested-r2} & Ratio of noisy affine and quadratic fit statistics\newline{\footnotesize\itshape proof} & \hyperref[gaussian:nested_r2]{p.~\pageref*{gaussian:nested_r2}}; \hyperref[gaussian:nested_ratio]{p.~\pageref*{gaussian:nested_ratio}}\\[4pt]

\codepath{eq:guide-scaling} & Proportional covariance scaling of the two components\newline{\footnotesize\itshape proof} & \hyperref[geometry:component_scaling]{p.~\pageref*{geometry:component_scaling}}; \hyperref[geometry:scaling_proposition]{p.~\pageref*{geometry:scaling_proposition}}\\[4pt]

\codepath{eq:guide-selected-moments} & Scalar selected-parent variance and mean shift\newline{\footnotesize\itshape proof} & \hyperref[selection:scalar_moments]{p.~\pageref*{selection:scalar_moments}}; \hyperref[selection:selected_moments]{p.~\pageref*{selection:selected_moments}}\\[4pt]

\codepath{eq:guide-response-bound-one} & One-dimensional curvature bound on relative displacement error\newline{\footnotesize\itshape proof + calculation} & \hyperref[selection:scalar_relative_error]{p.~\pageref*{selection:scalar_relative_error}}; \hyperref[selection:full_bound]{p.~\pageref*{selection:full_bound}}\\[4pt]

\codepath{eq:guide-recursion} & Scalar fixed-baseline inheritance closure\newline{\footnotesize\itshape definition + proof} & \hyperref[inheritance:scalar-map]{p.~\pageref*{inheritance:scalar-map}}; \hyperref[inheritance:retention-assumption]{p.~\pageref*{inheritance:retention-assumption}}\\[4pt]

\codepath{eq:guide-scalar-bounds} & Numerical endpoints of the scalar invariant interval\newline{\footnotesize\itshape calculation + proof} & \hyperref[inheritance:fixed-bounds]{p.~\pageref*{inheritance:fixed-bounds}}; \hyperref[inheritance:scalar-boundaries]{p.~\pageref*{inheritance:scalar-boundaries}}\\[4pt]

\codepath{eq:guide-general-fixedpoint-quadratic} & Polynomial satisfied by the scalar equilibrium\newline{\footnotesize\itshape proof} & \hyperref[inheritance:scalar-quadratic]{p.~\pageref*{inheritance:scalar-quadratic}}\\[4pt]

\codepath{eq:guide-scalar-root} & Positive scalar equilibrium root\newline{\footnotesize\itshape proof} & \hyperref[inheritance:positive-root]{p.~\pageref*{inheritance:positive-root}}; \hyperref[inheritance:rationalised-root]{p.~\pageref*{inheritance:rationalised-root}}\\[4pt]

\codepath{eq:guide-contraction} & Exact distance formula and scalar contraction inequality\newline{\footnotesize\itshape proof} & \hyperref[inheritance:contraction]{p.~\pageref*{inheritance:contraction}}; \hyperref[inheritance:contraction-bound]{p.~\pageref*{inheritance:contraction-bound}}\\[4pt]

\codepath{eq:guide-rho-one} & Scalar precision after repeated updates with full retention\newline{\footnotesize\itshape proof} & \hyperref[inheritance:precision-iteration]{p.~\pageref*{inheritance:precision-iteration}}; \hyperref[inheritance:rho-one]{p.~\pageref*{inheritance:rho-one}}\\[4pt]

\codepath{eq:guide-LA} & Scalar phenotype regression of relative fitness\newline{\footnotesize\itshape definition} & \hyperref[conversion:regression_model]{p.~\pageref*{conversion:regression_model}}; \hyperref[conversion:coefficient_convention]{p.~\pageref*{conversion:coefficient_convention}}\\[4pt]

\codepath{eq:guide-intercept} & Relative-fitness intercept and its numerical value\newline{\footnotesize\itshape proof + calculation} & \hyperref[conversion:phenotype_moment_formula]{p.~\pageref*{conversion:phenotype_moment_formula}}; \hyperref[conversion:scalar_example]{p.~\pageref*{conversion:scalar_example}}\\[4pt]

\codepath{eq:guide-first-moment} & First selected phenotype moment\newline{\footnotesize\itshape proof} & \hyperref[conversion:regression_moments]{p.~\pageref*{conversion:regression_moments}}; \hyperref[conversion:phenotype_moment_formula]{p.~\pageref*{conversion:phenotype_moment_formula}}\\[4pt]

\codepath{eq:guide-second-moment} & Second selected phenotype moment\newline{\footnotesize\itshape proof} & \hyperref[conversion:regression_moments]{p.~\pageref*{conversion:regression_moments}}; \hyperref[conversion:phenotype_moment_formula]{p.~\pageref*{conversion:phenotype_moment_formula}}\\[4pt]

\codepath{eq:guide-conditional-A-one} & Breeding-value mean and variance conditional on phenotype\newline{\footnotesize\itshape proof} & \hyperref[conversion:conditional_moments]{p.~\pageref*{conversion:conditional_moments}}; \hyperref[conversion:conditional_formula]{p.~\pageref*{conversion:conditional_formula}}\\[4pt]

\codepath{eq:guide-BV-moments} & Selected breeding-value moments in one trait\newline{\footnotesize\itshape proof} & \hyperref[conversion:selected_formula]{p.~\pageref*{conversion:selected_formula}}; \hyperref[conversion:scalar_example]{p.~\pageref*{conversion:scalar_example}}\\[4pt]

\codepath{eq:guide-conversion-one} & Scalar inversion from selected moments to log-quadratic parameters\newline{\footnotesize\itshape proof + calculation} & \hyperref[conversion:conversion_formula]{p.~\pageref*{conversion:conversion_formula}}; \hyperref[conversion:scalar_numbers]{p.~\pageref*{conversion:scalar_numbers}}\\[4pt]

\codepath{eq:guide-ebv} & Share under fixed-coefficient prediction-variance substitution\newline{\footnotesize\itshape proof + calculation} & \hyperref[estimation:shrinkage_share]{p.~\pageref*{estimation:shrinkage_share}}; \hyperref[estimation:shrinkage]{p.~\pageref*{estimation:shrinkage}}\\[4pt]

\codepath{eq:guide-trace-reduction} & Four-index reduction of quadratic variance to a trace\newline{\footnotesize\itshape proof} & \hyperref[gaussian:four_index_variance]{p.~\pageref*{gaussian:four_index_variance}}; \hyperref[gaussian:quadratic_variance]{p.~\pageref*{gaussian:quadratic_variance}}\\[4pt]

\codepath{eq:guide-multi-variances} & Multivariate linear variance, quadratic variance and orthogonality\newline{\footnotesize\itshape proof} & \hyperref[gaussian:variance_formulas]{p.~\pageref*{gaussian:variance_formulas}}; \hyperref[gaussian:orthogonality]{p.~\pageref*{gaussian:orthogonality}}\\[4pt]

\codepath{eq:guide-multi-selected} & Multivariate selected-parent mean and covariance\newline{\footnotesize\itshape proof} & \hyperref[selection:selected_moments]{p.~\pageref*{selection:selected_moments}}; \hyperref[selection:moment_formulas]{p.~\pageref*{selection:moment_formulas}}\\[4pt]

\codepath{eq:guide-whitened-full} & Exact selected displacement in covariance-scaled coordinates\newline{\footnotesize\itshape proof} & \hyperref[selection:whitened_error_identity]{p.~\pageref*{selection:whitened_error_identity}}; \hyperref[selection:precision_factorisation]{p.~\pageref*{selection:precision_factorisation}}\\[4pt]

\codepath{eq:guide-whitened-error} & Displacement error in covariance-scaled coordinates\newline{\footnotesize\itshape proof} & \hyperref[selection:whitened_error_identity]{p.~\pageref*{selection:whitened_error_identity}}\\[4pt]

\codepath{eq:guide-multi-bound} & Relative displacement bound and Frobenius upper bound\newline{\footnotesize\itshape proof} & \hyperref[selection:relative_bound]{p.~\pageref*{selection:relative_bound}}; \hyperref[selection:full_bound]{p.~\pageref*{selection:full_bound}}\\[4pt]

\codepath{eq:guide-genic-baseline} & Multitrait genic covariance from loci\newline{\footnotesize\itshape proof} & \hyperref[inheritance:genic-formula]{p.~\pageref*{inheritance:genic-formula}}; \hyperref[inheritance:genic]{p.~\pageref*{inheritance:genic}}\\[4pt]

\codepath{eq:guide-matrix-interval-map} & Invariance of the covariance order interval\newline{\footnotesize\itshape proof} & \hyperref[inheritance:invariance]{p.~\pageref*{inheritance:invariance}}; \hyperref[inheritance:existence]{p.~\pageref*{inheritance:existence}}\\[4pt]

\codepath{eq:guide-multi-conditioning} & Multivariate breeding-value conditioning on phenotype\newline{\footnotesize\itshape proof} & \hyperref[conversion:conditional_moments]{p.~\pageref*{conversion:conditional_moments}}; \hyperref[conversion:conditional_second]{p.~\pageref*{conversion:conditional_second}}\\[4pt]

\codepath{eq:guide-multi-raw-derivatives} & Derivatives of the logged phenotype-regression polynomial\newline{\footnotesize\itshape proof} & \hyperref[conversion:raw_derivatives]{p.~\pageref*{conversion:raw_derivatives}}; \hyperref[conversion:raw_log]{p.~\pageref*{conversion:raw_log}}\\[4pt]

\bottomrule\end{longtable}\normalsize

\Needspace{8\baselineskip}

\subsection{Supplementary File 3: worked estimation}

\small\begin{longtable}{@{}>{\raggedright\arraybackslash}p{0.28\textwidth}>{\raggedright\arraybackslash}p{0.47\textwidth}>{\raggedright\arraybackslash}p{0.16\textwidth}@{}}

\toprule Source label & Concept and treatment & Derivation\\\midrule\endhead

\codepath{eq:A} & Gaussian latent breeding-value population in the estimation example\newline{\footnotesize\itshape definition} & \hyperref[estimation:model]{p.~\pageref*{estimation:model}}; \hyperref[estimation:measurement_model]{p.~\pageref*{estimation:measurement_model}}\\[4pt]

\codepath{eq:measurement} & Repeated noisy trait measurements\newline{\footnotesize\itshape definition} & \hyperref[estimation:measurement_model]{p.~\pageref*{estimation:measurement_model}}; \hyperref[estimation:model]{p.~\pageref*{estimation:model}}\\[4pt]

\codepath{eq:count} & Conditional Poisson reproductive-count model\newline{\footnotesize\itshape definition} & \hyperref[estimation:count_model]{p.~\pageref*{estimation:count_model}}; \hyperref[estimation:model]{p.~\pageref*{estimation:model}}\\[4pt]

\codepath{eq:surface} & Quadratic log conditional expected count\newline{\footnotesize\itshape definition} & \hyperref[estimation:count_model]{p.~\pageref*{estimation:count_model}}\\[4pt]

\codepath{eq:linear} & Generating linear variance in the worked example\newline{\footnotesize\itshape calculation + proof} & \hyperref[estimation:components]{p.~\pageref*{estimation:components}}; \hyperref[gaussian:linear_variance]{p.~\pageref*{gaussian:linear_variance}}\\[4pt]

\codepath{eq:quadratic} & Generating quadratic variance in the worked example\newline{\footnotesize\itshape calculation + proof} & \hyperref[estimation:components]{p.~\pageref*{estimation:components}}; \hyperref[gaussian:quadratic_variance]{p.~\pageref*{gaussian:quadratic_variance}}\\[4pt]

\codepath{eq:share} & Generating variance share in the worked example\newline{\footnotesize\itshape calculation} & \hyperref[estimation:components]{p.~\pageref*{estimation:components}}; \hyperref[worked:estimation]{p.~\pageref*{worked:estimation}}\\[4pt]

\codepath{eq:Ehat} & Unbiased within-genotype measurement-variance calibration\newline{\footnotesize\itshape proof} & \hyperref[estimation:anova_calibration]{p.~\pageref*{estimation:anova_calibration}}; \hyperref[estimation:calibration_estimators]{p.~\pageref*{estimation:calibration_estimators}}\\[4pt]

\codepath{eq:Ghat} & Unbiased between-genotype variance calibration\newline{\footnotesize\itshape proof} & \hyperref[estimation:anova_calibration]{p.~\pageref*{estimation:anova_calibration}}; \hyperref[estimation:calibration_estimators]{p.~\pageref*{estimation:calibration_estimators}}\\[4pt]

\codepath{eq:shrink} & Calibrated prediction reliability and centred conditional mean\newline{\footnotesize\itshape proof + calculation} & \hyperref[estimation:posterior]{p.~\pageref*{estimation:posterior}}; \hyperref[estimation:reliability_identities]{p.~\pageref*{estimation:reliability_identities}}\\[4pt]

\codepath{eq:posterior} & Normal latent distribution conditional on repeated measurements\newline{\footnotesize\itshape proof} & \hyperref[estimation:posterior_formula]{p.~\pageref*{estimation:posterior_formula}}; \hyperref[estimation:posterior]{p.~\pageref*{estimation:posterior}}\\[4pt]

\codepath{eq:likelihood} & Poisson likelihood integrated over uncertain breeding value\newline{\footnotesize\itshape proof} & \hyperref[estimation:mixture_likelihood]{p.~\pageref*{estimation:mixture_likelihood}}; \hyperref[estimation:likelihood]{p.~\pageref*{estimation:likelihood}}\\[4pt]

\codepath{eq:derivatives} & Slope and curvature of the fitted quadratic log mean\newline{\footnotesize\itshape proof + calculation} & \hyperref[estimation:components]{p.~\pageref*{estimation:components}}\\[4pt]

\codepath{eq:quarticvar} & Complete even-part variance with a quartic term\newline{\footnotesize\itshape proof} & \hyperref[estimation:quartic_variance]{p.~\pageref*{estimation:quartic_variance}}; \hyperref[estimation:quartic_full_r2]{p.~\pageref*{estimation:quartic_full_r2}}\\[4pt]

\codepath{eq:quarticR2} & Full-curve affine R-squared for the quartic generating function\newline{\footnotesize\itshape proof + calculation} & \hyperref[estimation:quartic_full_r2]{p.~\pageref*{estimation:quartic_full_r2}}; \hyperref[estimation:quartic]{p.~\pageref*{estimation:quartic}}\\[4pt]

\bottomrule\end{longtable}\normalsize

\subsection{Unlabelled intermediate displays}

A further 87 unlabelled displays were inspected. They supply definitions, intermediate algebra and numerical substitutions within the derivations above. The machine-readable equation map groups them by concept and identifies their companion targets; no additional result was left without a derivation or calculation.

\section{The eleven interpretation questions}\label{coverage:questions}

The questions in Appendix S6 can be followed through the derivations as follows.

\begin{description}[style=nextline]

\item[Q1: what the share measures] A fraction of variance in a quadratic log conditional-mean description under the stated population. See \hyperref[gaussian:share]{p.~\pageref*{gaussian:share}}.

\item[Q2: established tools and the question addressed] The prerequisite sources and their roles are separated from the comparisons made here. See \hyperref[concepts:chapter]{Chapter~\ref*{concepts:chapter}}.

\item[Q3: why average before taking logs] Log conditional mean and conditional mean log are different targets. See \hyperref[conditioning:jensen]{p.~\pageref*{conditioning:jensen}}.

\item[Q4: gene action and a cutoff] A population fitness description alone identifies neither gene action nor a universal threshold. See \hyperref[gaussian:share]{p.~\pageref*{gaussian:share}}.

\item[Q5: which Gaussian assumptions matter] Third and fourth moments show which identities fail outside the Gaussian model. See \hyperref[gaussian:non_gaussian]{p.~\pageref*{gaussian:non_gaussian}}.

\item[Q6: estimation] Observation and latent-value uncertainty enter the likelihood and refitting procedure. See \hyperref[estimation:chapter]{Chapter~\ref*{estimation:chapter}}.

\item[Q7: zero gradient and zero variance] Zero shares, undefined ratios and reduced-rank supports are treated separately. See \hyperref[gaussian:singular_support]{p.~\pageref*{gaussian:singular_support}}.

\item[Q8: rankings and selection accuracy] Exact weighting and counterexamples distinguish a close fit from an accurate parental prediction. See \hyperref[selection:chapter]{Chapter~\ref*{selection:chapter}}.

\item[Q9: covariance, dimensions and orientation] Whitening, Rayleigh quotients and physical rotation answer different comparison questions. See \hyperref[geometry:chapter]{Chapter~\ref*{geometry:chapter}}.

\item[Q10: inheritance and equilibrium] The fixed baseline and retention rule are explicit modelling assumptions. See \hyperref[inheritance:chapter]{Chapter~\ref*{inheritance:chapter}}.

\item[Q11: use of the descriptive share] Report the share with both absolute components to distinguish different population changes. See \hyperref[worked:comparison]{p.~\pageref*{worked:comparison}}.

\end{description}

\section{Commands associated with the results}\label{coverage:commands}

Commands start at the repository root. Inputs, outputs and assumptions remain in \repofile{reference/result_catalog.json}{the machine-readable catalogue}. Executing these implementations does not replace the analytical arguments.

\begin{description}[style=nextline]

\item[main figures] Generate the four main figures from recorded parameters and seeds.
\begin{quote}\small\ttfamily\raggedright\detokenize{python code/figures.py}\end{quote}

\item[core mathematics] Check 30 matrix cases, coordinate changes, proportional scaling, bounds and scalar trajectories.
\begin{quote}\small\ttfamily\raggedright\detokenize{python code/verify_mathematics.py}\end{quote}

\item[result checks] Independent small numerical examples of the identities, scope conditions and counterexamples; no random simulation or document rebuild.
\begin{quote}\small\ttfamily\raggedright\detokenize{python code/reproduce_result_checks.py}\end{quote}

\item[simulation archive] Validate and summarise the 20 archived diagnostic tables; does not run SLiM.
\begin{quote}\small\ttfamily\raggedright\detokenize{python code/summarize_simulation.py}\end{quote}

\item[fresh simulation] Run seeds 42--61 with SLiM 5.2, reconstruct diagnostics and compare every diagnostic table with the archive.
\begin{quote}\small\ttfamily\raggedright\detokenize{bash run_all.sh --simulate}\end{quote}

\item[guide figures] Regenerate all five mathematics-guide figures and check their stated values.
\begin{quote}\small\ttfamily\raggedright\detokenize{python code/guide_figures.py}\end{quote}

\item[intro checks] Print quadrature checks of one-trait population regressions and selected-parent weighting.
\begin{quote}\small\ttfamily\raggedright\detokenize{python tutorials/introduction/check_examples.py}\end{quote}

\item[intro figures] Rebuild three introductory teaching figures from their specified examples.
\begin{quote}\small\ttfamily\raggedright\detokenize{python tutorials/introduction/make_figures.py}\end{quote}

\item[one trait] Calculate and quadrature-check the three one-trait populations, then draw Figure A1.
\begin{quote}\small\ttfamily\raggedright\detokenize{python tutorials/one_trait_visual/make_figure.py}\end{quote}

\item[geometry values] Quadrature-check all rotation and comparison cases without rebuilding visual documents.
\begin{quote}\small\ttfamily\raggedright\detokenize{python tutorials/geometry/code/geometry.py}\end{quote}

\item[geometry visuals] Rebuild the two-page geometry figure and offline explorer; Figure A2 is page one.
\begin{quote}\small\ttfamily\raggedright\detokenize{bash tutorials/geometry/run.sh}\end{quote}

\item[parental weighting] Recover the fixed parental cloud, compare frequency and reproductive weights, and rebuild its figures and offline explorer.
\begin{quote}\small\ttfamily\raggedright\detokenize{bash tutorials/parental_weighting/run.sh}\end{quote}

\item[estimation] Recreate the default observations, latent-value fit, 500 bootstrap refits and three estimation figures. The seed is 20261005; the bootstrap seed is 20261006.
\begin{quote}\small\ttfamily\raggedright\detokenize{bash run_estimation.sh}\end{quote}

\item[estimation checks] Run eight scientific checks: moments, observation/truth separation, independent integration, recovery, bootstrap calibration, quartic scope, integration failure and failed calibration.
\begin{quote}\small\ttfamily\raggedright\detokenize{python -m unittest discover -s examples/estimating_variance_share -p test_model.py}\end{quote}

\item[quartic case] Generate the optional misspecified quadratic analysis of data from a quartic log mean. These fitted results are an exploratory case; the cited 0.2424 is the exact generating full-curve R\textsuperscript{2}.
\begin{quote}\small\ttfamily\raggedright\detokenize{bash run_estimation.sh --quartic 0.08 --output outputs/estimation_quartic}\end{quote}

\item[shiny] Launch the local Shiny for Python teaching app. Its download records the displayed data, settings, estimates and any completed bootstrap.
\begin{quote}\small\ttfamily\raggedright\detokenize{bash run_shiny.sh}\end{quote}

\item[reference comparison] After a complete default run, compare figure parameters/results, archived-simulation summary, observations and every bootstrap draw with the reference records.
\begin{quote}\small\ttfamily\raggedright\detokenize{python code/check_results.py default}\end{quote}

\end{description}
